\begin{examplebox}
  \textbf{\textcolor{CExample}{例1（基础：由一个点确定幂函数）}}

  已知幂函数 $f(x)=x^\alpha$ 的图像经过 $(4,2)$。求指数，判断正半轴上的走势，并比较 $f(0.8)$ 与 $f(0.5)$。

  \textbf{解：} 代入得 $4^\alpha=2$，即 $2^{2\alpha}=2^1$，所以 $\alpha=1/2$。函数为 $f(x)=\sqrt{x}$，在 $x>0$ 时不断上升，但越来越平缓；其定义域还包括 $0$，所以完整图像经过原点。因为 $0.8>0.5$，故 $f(0.8)>f(0.5)$。

  \textbf{关键：} 点 $(4,2)$ 确定指数；指数确定后，就能判断走势，不必计算两个根号的具体值。

  \medskip
  \textbf{\textcolor{CExample}{例2（同底幂比较：不用计算小数次幂）}}

  比较 $1.7^{0.3}$ 与 $1.7^{0.5}$，以及 $0.8^{0.2}$ 与 $0.8^{0.3}$。

  \textbf{解：} 第一组底数大于 $1$，指数大的数更大，故 $1.7^{0.3}<1.7^{0.5}$；第二组底数小于 $1$ 但大于 $0$，指数大的数更小，故 $0.8^{0.2}>0.8^{0.3}$。

  \textbf{关键：} 一定先看底数与 $1$ 的关系，再看指数。

  \medskip
  \textbf{\textcolor{CExample}{例3（同指数比较：负指数反向）}}

  已知 $y=x^\alpha$ 在 $(0,+\infty)$ 上严格递减，且 $a^\alpha>b^\alpha$（$a,b>0$），求 $a,b$ 的大小关系。

  \textbf{解：} 由严格递减知 $\alpha<0$。若 $a\ge b$，利用递减性质就有 $a^\alpha\le b^\alpha$，与题设矛盾。故 $a<b$。

  \textbf{关键：} 负指数等价于取正指数幂的倒数，比较方向反过来。

  \medskip
  \textbf{\textcolor{CExample}{例4（定义应用：先令系数等于 $1$）}}

  已知 $f(x)=(9m^2-3)x^m$ 是幂函数，图像经过点 $(n,-1/m)$。判断下列说法：
  \begin{enumerate}[label=\Alph*.,leftmargin=*,itemsep=0.25em]
    \item $m=-\dfrac23$；
    \item $f(x)$ 是偶函数；
    \item $n=\dfrac{3\sqrt6}{4}$；
    \item $f(a+1)>f(3-a)$ 的解集为 $(-\infty,1)$。
  \end{enumerate}

  \textbf{解：} 这是用户所提供题图中的多选题。首先，由幂函数的定义（或必过 $(1,1)$）得到
  \[
     9m^2-3=1,\qquad m=\pm\frac23.
  \]
  若 $m=2/3$，则 $f(x)=(\sqrt[3]{x})^2\ge0$，不可能经过纵坐标为 $-3/2$ 的点。故 $m=-2/3$，A 正确。

  此时按实数范围内奇次根的意义，
  \[
    f(x)=\frac1{(\sqrt[3]{x})^2}
        =\frac1{|x|^{2/3}},\qquad x\ne0.
  \]
  定义域关于原点对称，且 $f(-x)=f(x)$，B 正确。代入题设点：
  \[
  \frac1{|n|^{2/3}}=\frac32
  \quad\Longrightarrow\quad
  |n|=\left(\frac23\right)^{3/2}
     =\frac{2\sqrt6}{9}.
  \]
  因此 $n=\pm 2\sqrt6/9$，C 错误。

  判断 D 时，不能把整个定义域当作递减区间。先排除 $a=-1,3$，再利用偶函数的性质和正半轴上的递减性：
  \[
  \begin{aligned}
  f(a+1)>f(3-a)
  &\iff |a+1|<|3-a|\\
  &\iff (a+1)^2<(3-a)^2\\
  &\iff 8a<8\\
  &\iff a<1.
  \end{aligned}
  \]
  结合定义域，解集为 $(-\infty,-1)\cup(-1,1)$。D 漏掉了 $a=-1$ 不存在函数值这一点，故错误。

  \textbf{答案：A、B。} 本题最关键的入口是系数等于 $1$；后续则用到了幂函数的值域、奇偶性、单调性和定义域。特别注意：$x^{\pm2/3}$ 按奇次根解释可以有负数底数，不能把正半轴上的规律直接套到负半轴。

  \medskip
  \textbf{\textcolor{CExample}{例5（图像交点与不等式）}}

  解不等式 $x^{0.3}<x^{0.5}$，限定 $x>0$。

  \textbf{解：} 指数 $0.3<0.5$。当 $x>1$ 时成立；当 $0<x<1$ 时恰好相反；$x=1$ 时相等，不满足严格不等式。因此解集为 $(1,+\infty)$。

  也可以两边除以正数 $x^{0.3}$，得到 $1<x^{0.2}$，即 $x>1$。若把两条曲线画出来，这正是它们在 $(1,1)$ 右侧的上下关系。

  \medskip
  \textbf{\textcolor{CExample}{例6（推广：改变指数与改变系数是两回事）}}

  已知 $g(x)=Ax^{-2}$ 经过 $(1,3)$，求 $A$；它是否是高中定义下的幂函数？再判断 $g(x)$ 在 $x>0$ 时的走势。

  \textbf{解：} $g(1)=A=3$，所以 $g(x)=3/x^2$。它不是严格定义下的幂函数，但可由 $x^{-2}$ 的图像在竖直方向放大三倍得到，仍然在正半轴上递减。若系数改成 $-3$，图像关于 $x$ 轴翻转，在正半轴上将变为递增。

  \textbf{关键：} 同一个指数，系数的正负也能影响增减方向。这正是研究基本幂函数时把系数固定为 $1$ 的好处。

  \medskip
  \textbf{\textcolor{CExample}{例7（顿悟：不同底数、不同指数时，用 $1$ 作桥梁）}}

  比较 $0.6^{0.7}$ 与 $1.2^{0.3}$。

  \textbf{解：} 由于两个指数均为正数，第一个底数小于 $1$，第二个底数大于 $1$，故
  \[
      0.6^{0.7}<1<1.2^{0.3}.
  \]
  不需要计算任何一个幂的具体数值。但如果两数都在 $1$ 的同一侧，仅靠这一条不能直接判定大小。
  \medskip
  \textbf{\textcolor{CExample}{例8（图像反推指数：公共点处的切线）}}

  若幂函数 $y=x^\alpha$ 在 $(1,1)$ 处的切线为 $y=-2x+3$，求该幂函数。

  \textbf{解：} 切线斜率为 $-2$。由于 $f'(1)=\alpha$，直接得到 $\alpha=-2$，故 $y=1/x^2$。无需求出其他点的坐标。
\end{examplebox}
